We have presented a comprehensive empirical study of hyperparameter
scaling laws in LLM pre-training, focusing on weight decay and batch
size. Our approach leverages the AdamW timescale ($\tepoch$) to
develop robust scaling relationships that predict optimal
hyperparameter settings across a broad spectrum of model ($N$),
dataset ($D$), and batch sizes ($B$). We demonstrated that optimal
$\tepoch$ decreases as a power law with the tokens-per-parameter
ratio, providing a systematic method to set weight decay optimally
across diverse training scenarios.

Furthermore, we introduced a novel, practical methodology for
estimating critical batch size ($\bcrit$). Our findings diverge from
influential prior work that tied $\bcrit$ predominantly to compute or
loss, while agreeing with the recent findings of~\citet{zhang2024how}
that underscore dataset size as the principal scaling factor.
Additionally, we showed that contrary to previous studies suggesting
optimal batch size ($\bopt$) scales primarily with compute, it also
exhibits a clear power-law dependence on $D$.

Also, our analysis of Pareto-optimal configurations reveals an
important strategic advantage for smaller, overtrained models in
scenarios where rapid training and high parallelism are prioritized.

\cref{sec:limitations} notes limitations and directions for
further study suggested by our results.
%
In particular, as inference-time scaling comes to the fore, inference
time and compute must also be considered as first-class Pareto
objectives.  Moreover, finer-grained configuration decisions, such as
model depth and context length, should be considered along with $N$,
$D$, and $B$.
